\chapter{Computational Validation and Numerical Regression}
\label{chap:computation}

\section{Role of computation}

Computation serves three legitimate purposes in this project:
\begin{enumerate}
  \item verify that symbolic formulas were implemented correctly;
  \item visualize geometry and local stability;
  \item test candidate constructions for immediate counterexamples.
\end{enumerate}
It does not replace proof of a universal statement. A finite list of critical-line zeros is compatible with both RH and the possibility of a later off-axis zero.

The scripts supplied with this monograph generate every figure and numerical table from source. No zero locations are hard-coded. The principal dependencies are Python, NumPy, Matplotlib, and mpmath.

\section{Elementary unit tests}

The repository's initial code implements:
\begin{align*}
 \Hbin(\sigma)&=-\sigma\ln\sigma-(1-\sigma)\ln(1-\sigma),\\
 \Acomp(\sigma,\phi)&=\sqrt\sigma+e^{i\phi}\sqrt{1-\sigma},\\
 \J(s)&=1-\overline{s}.
\end{align*}
The baseline unit tests check
\begin{align*}
 \Hbin(1/2)&=\ln2,\\
 \Acomp(1/2,\pi)&=0,\\
 \J(1/2+17i)&=1/2+17i.
\end{align*}
The monograph adds property-style tests across sampled domains: entropy symmetry and strict deficit, involution idempotence, cancellation uniqueness to numerical tolerance, eta--zeta identity, and xi functional symmetry.

\section{Functional-equation residuals}

The completed function is evaluated at high precision using
\[
 \xi(s)=\frac12s(s-1)\pi^{-s/2}\Gamma(s/2)\zeta(s).
\]
At selected points, the residuals
\[
 |\xi(s)-\xi(1-s)|,
 \qquad
 |\xi(\J s)-\overline{\xi(s)}|
\]
are at the scale expected from the chosen arbitrary precision. These checks confirm code consistency; they do not prove the functional equation, which was derived analytically in \cref{chap:zeta}.

\begin{table}[htbp]
\centering
\caption{High-precision residual checks for the functional equation and conjugate-affine involution.}
\label{tab:fe-residuals}
\begin{tabular}{rrrr}
\toprule
$\sigma$ & $t$ & $|\xi(s)-\xi(1-s)|$ & $|\xi(Js)-\overline{\xi(s)}|$\\
\midrule
0.23 & 7.1 & 3.61263205064e-82 & 3.61263205064e-82\\
0.41 & 14.2 & 6.43444698684e-85 & 6.43444698684e-85\\
0.77 & 22.5 & 6.24309244359e-86 & 6.24309244359e-86\\
0.5 & 33.0 & 2.71847215114e-90 & 2.71847215114e-90\\
1.3 & 4.0 & 1.97666211436e-82 & 1.97666211436e-82\\
\bottomrule
\end{tabular}
\end{table}


\section{Critical-line visualization}

On the critical line, $\XiR(t)=\xi(1/2+it)$ is real. \Cref{fig:xi-line} plots a scaled version for $0\le t\le50$. Vertical markers indicate the first zeros returned by mpmath.

\begin{figure}[htbp]
 \centering
 \includegraphics[width=0.86\textwidth]{xi_critical_line.png}
 \caption{A scaled numerical plot of $\XiR(t)$ on the critical line. The finite plot is illustrative and not a proof of global zero location.}
 \label{fig:xi-line}
\end{figure}

The plot emphasizes the separation between the balance coordinate and spectral coordinate. Restricting to $\sigma=1/2$ makes the completed function real; zeros then occur at discrete ordinates. The ansatz seeks a mechanism that proves this restriction is compulsory rather than chosen for plotting.

\section{Cancellation landscape}

The squared amplitude
\[
 1+2\sqrt{\sigma(1-\sigma)}\cos\phi
\]
is evaluated on a grid. Its unique zero in one phase period is visible in \cref{fig:cancellation-landscape}. The local approximation
\[
 2x^2+\frac12\delta^2
\]
is compared against the exact expression in \cref{fig:cancellation-quadratic}. These plots are direct visual checks of the proof, not empirical discoveries.

\section{Eta prefactor test}

The factor $1-2^{1-s}$ is evaluated across a region containing the critical strip. The numerical valleys align with $\Re s=1$ at the predicted ordinates $-2\pi k/\ln2$, as shown in \cref{fig:eta-prefactor}. The factor remains bounded away from zero on compact subsets of the open strip, consistent with \cref{prop:eta-zeta-equivalence}.

A regression test samples points with $0<\sigma<1$ and verifies
\[
 \eta(s)-(1-2^{1-s})\zeta(s)
\]
to high precision. Near $s=1$, direct evaluation requires care because the factor and pole cancel; mpmath's analytic continuation handles the limit, while a dedicated series expansion is preferable for rigorous interval computation.

\section{First-zero table}

The first twenty non-trivial zeros are included in \cref{app:numerical}. Their real parts are reported as $0.5$ by the numerical routine. This table has two purposes: it fixes a reproducible software baseline and supplies test ordinates for future candidate kernels. It is not used to fit any parameter in the ansatz.

\section{Falsification experiments for candidate maps}

Any proposed state map should undergo at least the following tests before analytic effort is invested:
\begin{enumerate}
  \item \textbf{Zero coincidence.} Does the detector vanish at computed zeta zeros to increasing precision?
  \item \textbf{False positives.} Does it vanish at non-zeros on or off the critical line?
  \item \textbf{Covariance.} Is the state at $\J s$ related to the swapped conjugate state?
  \item \textbf{Metric stability.} Are the component weights in $[0,1]$ and do they sum to one?
  \item \textbf{Gauge stability.} Is the zero readout independent of arbitrary phase conventions?
  \item \textbf{Approach stability.} At a zero, does the state ray have a unique limit from different directions?
  \item \textbf{No fitting.} Does the algorithm work when evaluated at zeros not used during development?
  \item \textbf{Multiplicity.} Does local scaling match the expected order of the zero?
\end{enumerate}
A candidate failing any basic test should be rejected or reformulated before interpretive claims are made.

\section{Interval arithmetic and proof assistants}

Floating-point tests are not certified. If a candidate bridge reaches a serious stage, the next computational layer should include:
\begin{itemize}
  \item interval arithmetic for bounds on kernels and non-vanishing factors;
  \item exact symbolic differentiation for local expansions;
  \item formal verification of elementary lemmas in Lean, Isabelle, or Coq;
  \item reproducible containers with pinned dependency versions;
  \item hash manifests for generated figures and tables.
\end{itemize}
Formalizing the cancellation theorem is easy; formalizing the analytic continuation and spectral construction would be substantially harder but valuable.

\section{Reproducibility manifest}

The asset script writes figures and numerical tables into the repository build directories. The PDF build should proceed from a clean checkout by running:
\begin{verbatim}
python scripts/generate_monograph_assets.py
cd monograph
latexmk -pdf main.tex
\end{verbatim}
The generated PDF, log, auxiliary files, and rendered page checks should be recorded in an append-only build log before release.

\status{The numerical results are regression evidence only. Every theorem used in the logical chain is proved independently of the computations.}
